%%%%%%%%%%  scan idx 175  |  folio 162  %%%%%%%%%%
thèses (1.1 (i) et (ii)
respée). Notons $\mathcal C_{1S}$ la sous-catégorie strictement pleine de $\mathcal C_S$ formée des objets qui sont localement dans $\mathcal C_0$. Alors :
\begin{enumerate}[label=\alph*)]
\item $\mathcal C_{1S}$ est stable par noyau d'épimorphisme (I 1.2) et quasi-relevable dans $\mathcal C_S$ (I 1.2) (en tant que catégorie fibrée sur un site ponctuel).
\item Soient $E\in\Ob\mathrm D(\mathcal C_S)$ et $[m,n]$ un intervalle de $\mathbb Z$. Pour que $E$ soit de $\mathcal C_0$-amplitude parfaite contenue dans $[m,n]$, il faut et il suffit que $E$ soit isomorphe, dans $\mathrm D(\mathcal C_S)$, à un complexe $F$ tel que $F^i=0$ pour $i\notin[m,n]$, $F^i\in\Ob\mathcal C_{0S}$ pour $i>m$, et $F^m\in\Ob\mathcal C_{1S}$.
\item Le foncteur canonique
\[
\mathrm K^b(\mathcal C_{1S})/\mathrm K^{b,\varnothing}(\mathcal C_{1S})
\longrightarrow \mathrm D(\mathcal C_S)
\]
est pleinement fidèle et a pour image essentielle la sous-catégorie pleine de $\mathrm D(\mathcal C_S)$ formée des complexes d'amplitude $\mathcal C_0$-parfaite finie.
\item Si $\mathcal C_{1S}$ est non seulement quasi-relevable mais relevable dans $\mathcal C_S$, on peut remplacer dans b) l'expression « $E$ soit isomorphe à $F$ dans $\mathrm D(\mathcal C_S)$ » par « il existe un quasi-isomorphisme $F\to E$ », et dans c) la catégorie $\mathrm K^{b,\varnothing}(\mathcal C_{1S})$ est réduite aux objets nuls.
\end{enumerate}

\emph{Preuve.} a) Il est clair que $\mathcal C_{1S}$ est stable par noyau d'épimorphisme. D'autre part, les objets de $\mathcal C_{1S}$ sont $\mathcal C_0$-pseudo-cohérents donc $\mathcal C_{0S}$-pseudo-cohérents d'après (1.1 c)) ; la deuxième assertion résulte donc de (I 2.14) (où l'on prend pour $S$ le site ponctuel).

b) Il suffit de démontrer la nécessité. On sait déjà, grâce à (1.1 c)), que $E$ est $\mathcal C_{0S}$-pseudo-cohérent. Comme $E$ est acyclique en degré $>n$, on peut donc supposer (I 2.10), quitte à remplacer $E$ par un complexe isomorphe dans $\mathrm D^{-}(\mathcal C_S)$, que $E$ est strictement $(m+1)$-$\mathcal C_{0S}$-pseudo-cohérent et que $E^i=0$ pour $i>n$. Comme $E$ est acyclique en degré $<m$, on peut en outre supposer, quitte à remplacer $E$ par son tronqué
\[
0\longrightarrow E^m/B^m\longrightarrow E^{m+1}\longrightarrow\ldots\longrightarrow E^n\longrightarrow0,
\]
que $E^i=0$ pour $i<m$. Alors, d'après (I 4.16), $E^m$ est localement objet de $\mathcal C_0$, donc objet de $\mathcal C_{1S}$, ce qui démontre l'assertion.

